\begin{proofbox}
	\textbf{\textcolor{CProof}{思路提示}}
	三次函数是多项式函数，对称中心必然满足恒等式$f(x_0+t)+f(x_0-t)=2f(x_0)$。我们可以先通过二阶导找到可能的拐点横坐标，再验证该点满足对称条件。

	\medskip
	\textbf{\textcolor{CProof}{正式推导}}
	\begin{description}[style=nextline, leftmargin=2.8em, labelsep=0.5em, itemsep=0.55em]
		\item[\textbf{步骤一（求拐点横坐标）：}]
			计算二阶导 $f''(x)=6ax+2b$。令 $f''(x)=0$，解得 $x_0=-\dfrac{b}{3a}$。
			\par\textbf{理由：} 三次函数有唯一拐点，拐点处二阶导为零。
		\item[\textbf{步骤二（计算纵坐标）：}]
			代入原函数得 $y_0=f(x_0)$，即 $y_0=ax_0^3+bx_0^2+cx_0+d$。
			\par\textbf{理由：} 对称中心必在函数图像上。
		\item[\textbf{步骤三（构造对称差式）：}]
			设 $\Delta(t)=f(x_0+t)+f(x_0-t)-2y_0$，目标证明 $\Delta(t)\equiv 0$。
			\par\textbf{理由：} 中心对称的充要条件。
		\item[\textbf{步骤四（展开合并）：}]
			分别展开 $f(x_0\pm t)$：
			\[
				\begin{aligned}
					f(x_0\pm t) &= a(x_0^3\pm 3x_0^2t+3x_0t^2\pm t^3) \\ &\quad + b(x_0^2\pm 2x_0t + t^2) + c(x_0\pm t) + d.
			\end{aligned}
			\]
			两式相加，奇次项抵消，偶次项加倍：
			\[
				f(x_0+t)+f(x_0-t) = 2(ax_0^3+bx_0^2+cx_0+d) + 2(3a x_0+b)t^2.
			\]
			\par\textbf{理由：} 直接代数运算。
		\item[\textbf{步骤五（代入化简）：}]
			将 $x_0=-\dfrac{b}{3a}$ 代入 $3a x_0+b$，得 $3a\left(-\frac{b}{3a}\right)+b=0$。于是
			\[
				f(x_0+t)+f(x_0-t)=2y_0.
			\]
			即 $\Delta(t)=0$ 对所有 $t$ 成立。
			\par\textbf{理由：} 利用第一步求出的 $x_0$ 化简系数。
	\end{description}

	\medskip
	\textbf{\textcolor{CProof}{结论回扣}}
	我们证明了三次函数 $f(x)=ax^3+bx^2+cx+d\ (a\neq0)$ 的图像关于点 $\left(-\dfrac{b}{3a},\ f\!\left(-\dfrac{b}{3a}\right)\right)$ 中心对称，且该点就是函数的拐点。这一结论可以直接用于求解对称中心及利用对称性计算函数值。
\end{proofbox}
